% ============================================================================
% fondos.tex  —  CONTEXTOS / escenarios planos para el capibara (obra propia, CC0)
% ----------------------------------------------------------------------------
% Fondos flat, limpios, tipo escenario. Se dibujan DETRAS del capibara.
% El capibara (base ~+-1.7 en x, pies en y~-1.96) va CENTRADO al frente; el piso
% de cada escena esta a y~-1.95 para que "pise" el suelo.
%
% Requiere: % ============================================================================
% _estilo_interes.tex  —  Estilo compartido del "mundo del capibara" (CC0)
% ----------------------------------------------------------------------------
% Paleta plana + estilo de linea + FRAME canonico para trajes/accesorios que se
% dibujan EN LAS MISMAS COORDENADAS que el capibara base
% (tikzlib/interes/capibara_kawaii.tex). Todo obra propia, CC0.
%
% Requiere: \usepackage{tikz}. NO usa texto/fuentes -> apto SVGMobject de Manim.
%
% FRAME CANONICO DE REGISTRO (\capiframe):
%   Todos los trajes/accesorios y el propio capibara se exportan con ESTE mismo
%   \useasboundingbox. Asi los SVG comparten viewBox y en Manim se pueden cargar
%   con  SVGMobject(..., should_center=False)  y quedan PERFECTAMENTE alineados
%   sobre el capibara sin tocar nada (registro automatico).
%   Con should_center=True (por defecto) el frame se ignora (Manim recorta al
%   contenido visible), asi que NO rompe usos existentes del capibara.
% ============================================================================

% --- linework del capibara (idempotente con capibara_kawaii.tex) ---
\providecolor{capicontorno}{HTML}{6E4E34}  % linea (marron calido suave)
\tikzset{
  capilinea/.style   ={draw=capicontorno, line width=2.1pt,
                       line cap=round, line join=round},
  capilineaf/.style  ={draw=capicontorno, line width=1.6pt,
                       line cap=round, line join=round}, % linea fina
}

% --- paleta de trajes/accesorios (plana, tierna) ---
\providecolor{vueltiaobase}{HTML}{F2E6C6}  % cana flecha (crema claro)
\providecolor{vueltiaosom} {HTML}{E4D2A6}  % sombra del ala
\providecolor{vueltiaopin} {HTML}{3B332B}  % "pintao" (bandas oscuras)

\providecolor{banqcorbata} {HTML}{9C3B4C}  % corbatin vino
\providecolor{banqcorbdark}{HTML}{7E2E3D}  % corbatin sombra
\providecolor{oro}         {HTML}{D2A63C}  % oro (monoculo/detalles)
\providecolor{orodark}     {HTML}{B4862A}

\providecolor{comerdelan}  {HTML}{5E86A6}  % delantal azul mercado
\providecolor{comerdeldark}{HTML}{4B6E8C}
\providecolor{comergorra}  {HTML}{CF5B45}  % gorra roja
\providecolor{comergordark}{HTML}{B24734}

\providecolor{acadboard}   {HTML}{343747}  % birrete azul-noche
\providecolor{acadbdark}   {HTML}{262838}
\providecolor{acadgafas}   {HTML}{3B332B}  % marco de gafas

\providecolor{ahorrosa}    {HTML}{F2A6BC}  % cerdito rosado
\providecolor{ahorrodark}  {HTML}{E086A0}
\providecolor{ahorrogorro} {HTML}{D9A441}  % gorro mostaza

\providecolor{naranjafill} {HTML}{F59E2E}  % naranja
\providecolor{naranjadark} {HTML}{E07C1A}  % naranja sombra
\providecolor{naranjacont} {HTML}{9A5A1E}  % contorno naranja
\providecolor{hojaverde}   {HTML}{6BA83C}  % hoja
\providecolor{hojadark}    {HTML}{4E8A2E}
\providecolor{canastamimb} {HTML}{C8925A}  % mimbre canasta
\providecolor{canastadark} {HTML}{A9743F}

% --- fondo/escena ---
\providecolor{escpared}    {HTML}{FCF4E2}  % pared crema (combina con el video)
\providecolor{escpiso}     {HTML}{EFE1C4}  % piso
\providecolor{esczocalo}   {HTML}{E4D2A6}
\providecolor{banqcol}     {HTML}{D9E3EC}  % columnas banco
\providecolor{banqcoldark} {HTML}{BFCEDC}
\providecolor{banqarco}    {HTML}{C7D6E4}
\providecolor{merctoldoA}  {HTML}{D9694E}  % toldo rojo
\providecolor{merctoldoB}  {HTML}{FBF1DC}  % toldo crema
\providecolor{mercmadera}  {HTML}{C8925A}  % madera del puesto
\providecolor{mercmaddark} {HTML}{A9743F}
\providecolor{fincacielo}  {HTML}{CDE7F0}  % cielo
\providecolor{fincapasto}  {HTML}{9FCB6B}  % pasto
\providecolor{fincapastod} {HTML}{7FB050}
\providecolor{fincasol}    {HTML}{F6C64B}  % sol
\providecolor{fincamata}   {HTML}{6BA83C}

% FRAME CANONICO (capibara + trajes) — simetrico en x; contiene sombreros altos
% (hasta y~2.05) y la alcancia lateral (hasta x~2.65).
\newcommand{\capiframe}{\useasboundingbox (-2.75,-2.10) rectangle (2.75,2.25);}

% Uso:  \fondoneutro  \fondobanco  \fondomercado  \fondofinca
% Escala: dibujados para un lienzo 16:9 (\escenaframe). El capibara se coloca
%         encima a escala 1 (misma unidad). Recoloreables via _estilo_interes.tex.
% ============================================================================

% FRAME de escena 16:9 (para los fondos)
\newcommand{\escenaframe}{\useasboundingbox (-6.40,-3.60) rectangle (6.40,3.60);}

% piso comun (rectangulo inferior) — top del piso en y=-1.95
\newcommand{\escpisobase}[1]{% #1 = color del piso
  \path[fill=#1] (-6.40,-3.60) rectangle (6.40,-1.95);
}

% ---------------------------------------------------------------------------
% NEUTRO: pared crema + piso + halo suave detras del personaje. Muy limpio.
% ---------------------------------------------------------------------------
\newcommand{\fondoneutro}{%
  \begin{scope}
    \path[fill=escpared] (-6.40,-3.60) rectangle (6.40,3.60);
    \escpisobase{escpiso}
    % zocalo
    \path[fill=esczocalo] (-6.40,-1.95) rectangle (6.40,-1.70);
    % halo / spot detras del personaje
    \path[fill=escpiso, opacity=0.55] (0,-0.10) ellipse[x radius=2.7, y radius=2.7];
    % sombra del piso donde se para el personaje
    \path[fill=esczocalo] (0,-1.98) ellipse[x radius=1.7, y radius=0.30];
  \end{scope}
}

% ---------------------------------------------------------------------------
% BANCO: pared + piso + 2 columnas + arco/fronton central + mostrador lateral.
% ---------------------------------------------------------------------------
\newcommand{\fondobanco}{%
  \begin{scope}
    \path[fill=escpared] (-6.40,-3.60) rectangle (6.40,3.60);
    \escpisobase{escpiso}
    % fronton / arco central (semicirculo + cornisa)
    \path[fill=banqarco] (0,2.05) ellipse[x radius=3.05, y radius=1.05];
    \path[fill=banqcoldark] (-3.20,1.85) rectangle (3.20,2.15);
    % columnas (izq y der), con capitel y base
    \foreach \cx in {-3.55,3.55}{
      \path[fill=banqcol] (\cx-0.42,-1.95) rectangle (\cx+0.42,1.95);
      \path[fill=banqcoldark] (\cx-0.52,1.80) rectangle (\cx+0.52,2.05); % capitel
      \path[fill=banqcoldark] (\cx-0.52,-1.95) rectangle (\cx+0.52,-1.70); % base
      % estrias
      \draw[capilineaf, draw=banqcoldark] (\cx-0.18,-1.65) -- (\cx-0.18,1.72);
      \draw[capilineaf, draw=banqcoldark] (\cx+0.18,-1.65) -- (\cx+0.18,1.72);
    }
    % zocalo
    \path[fill=esczocalo] (-6.40,-1.95) rectangle (6.40,-1.75);
    % mostrador lateral (a la derecha, no tapa al personaje)
    \path[capilineaf, fill=mercmadera] (2.35,-1.95) rectangle (5.30,-0.65);
    \path[capilineaf, fill=mercmaddark] (2.20,-0.65) rectangle (5.45,-0.42);
    % sombra del piso
    \path[fill=esczocalo] (0,-1.98) ellipse[x radius=1.8, y radius=0.30];
  \end{scope}
}

% ---------------------------------------------------------------------------
% MERCADO: puesto con toldo a rayas rojo/crema + mesa de madera + postes.
% ---------------------------------------------------------------------------
\newcommand{\fondomercado}{%
  \begin{scope}
    \path[fill=escpared] (-6.40,-3.60) rectangle (6.40,3.60);
    \escpisobase{escpiso}
    % postes del puesto
    \path[fill=mercmaddark] (-4.60,-1.95) rectangle (-4.30,2.55);
    \path[fill=mercmaddark] ( 4.30,-1.95) rectangle ( 4.60,2.55);
    % viga superior
    \path[fill=mercmadera] (-4.75,2.35) rectangle (4.75,2.60);
    % toldo a rayas (festón ondulado abajo)
    \foreach \i [evaluate=\i as \x using -4.5+\i*1.0] in {0,...,9}{
      \pgfmathparse{Mod(\i,2)==0 ? "merctoldoA" : "merctoldoB"}\edef\tcol{\pgfmathresult}
      \path[fill=\tcol] (\x,2.35) rectangle (\x+1.0,1.55);
      \path[fill=\tcol] (\x+0.5,1.55) ellipse[x radius=0.5, y radius=0.28];
    }
    % mesa/tablero del puesto (delante, bajo, no tapa la cara)
    \path[capilineaf, fill=mercmadera] (-4.30,-1.95) rectangle (4.30,-0.95);
    \path[capilineaf, fill=mercmaddark] (-4.55,-0.95) rectangle (4.55,-0.68);
    % vetas de la madera
    \draw[capilineaf, draw=mercmaddark] (-4.0,-1.45) -- (4.0,-1.45);
    % zocalo
    \path[fill=esczocalo] (-6.40,-1.95) rectangle (6.40,-1.78);
  \end{scope}
}

% ---------------------------------------------------------------------------
% FINCA: cielo + sol + colinas + pasto + matas. Campo colombiano soleado.
% ---------------------------------------------------------------------------
\newcommand{\fondofinca}{%
  \begin{scope}
    % cielo
    \path[fill=fincacielo] (-6.40,-3.60) rectangle (6.40,3.60);
    % sol con rayos (arriba-derecha)
    \foreach \a in {0,30,...,330}
      \draw[line width=3pt, line cap=round, draw=fincasol]
          (4.30,2.35) -- +(\a:1.15);
    \path[fill=fincasol] (4.30,2.35) circle[radius=0.72];
    % colinas (dos arcos)
    \path[fill=fincapastod] (-6.40,-1.95) .. controls (-3.0,-0.35) and (0.0,-1.10) .. (6.40,-0.55) -- (6.40,-3.60) -- (-6.40,-3.60) -- cycle;
    % pasto (delante)
    \path[fill=fincapasto] (-6.40,-3.60) rectangle (6.40,-1.60);
    \draw[capilineaf, draw=fincapastod] (-6.40,-1.60) .. controls (0,-1.35) .. (6.40,-1.60);
    % matas / arbustos
    \foreach \bx in {-5.2,-3.4,3.2,5.0}{
      \path[fill=fincamata] (\bx,-1.65) circle[radius=0.42];
      \path[fill=fincamata] (\bx-0.35,-1.80) circle[radius=0.30];
      \path[fill=fincamata] (\bx+0.35,-1.80) circle[radius=0.30];
    }
    % nubes
    \path[fill=escpared] (-3.6,2.2) circle[radius=0.5];
    \path[fill=escpared] (-3.0,2.3) circle[radius=0.62];
    \path[fill=escpared] (-2.35,2.2) circle[radius=0.48];
    % sombra del personaje en el pasto
    \path[fill=fincapastod] (0,-1.95) ellipse[x radius=1.7, y radius=0.28];
  \end{scope}
}
